% ===========================================================================
\ReportSection[sec:req]{Problem Statement and Requirements}

\ReportSubsection{Operational needs}
\TodoAnchor{aoa-storage}

\textbf{Computer~1 (NAS and private cloud)}
\begin{itemize}
  \item Departmental file share (SMB, with NFS for the ML nodes), with access
    control lists, snapshots and quotas. Capacity target: \PH{N~TB now, growth
    \%/yr}.
  \item Bulk project and data storage; backup target for the ML nodes.
  \item License server hosting for \PH{named engineering software}.
  \item A private cloud for specific software: virtual machines and containers
    for \PH{named applications, e.g., version control, container registry,
    notebooks, wiki}.
  \item The LLM application tier: gateway and API, chat interface, retrieval
    service, vector database, model and artifact registry, centralized logging.
\end{itemize}

\textbf{Computer~2 (ML node[s])}
\begin{itemize}
  \item Serve an LLM to at least 20 named users, about 12 of them at the same
    time, with at least 128K-token context per request.
  \item Run fine-tuning and evaluation jobs (adapter-style fine-tuning of
    open-weight models, embedding and reranker adaptation, small custom ML
    models) \emph{at the same time} as serving, when two nodes are provisioned.
  \item Provide a backup path for serving if the primary node is unavailable.
  \item Provide an OpenAI-compatible API for scripts and pipelines, and
    retrieval over controlled repositories with per-group access control.
  \item Use the archived data on the NAS: answer questions from it with citations
    (retrieval), and train on selected documents, such as a licensed copy of NFPA~780,
    so the model handles that material well (Section~\ref{sec:archive}).
  \item Optionally load the same serving model on Node~B so it can act as a second
    serving node when it is not training (Section~\ref{sec:nodeb}).
\end{itemize}

\ReportSubsection{Compliance requirements (CUI)}
\TodoAnchor{aoa-cui}

\TableTitle{CUI-related requirements and their implications}
{\footnotesize
\begin{xltabular}{\textwidth}{@{}l Y{0.88} Y{1.12}@{}}
  \toprule
  Area & Driver & Implication \\
  \midrule
  \endhead
  Safeguarding & NIST SP 800-171 (Rev~2 or 3 per contract); 32~CFR Part~2002 & Access control, audit, encryption and boundary protection for every component, including the NAS. \\
  Contractual & DFARS 252.204-7012; CMMC Level~2 (if applicable) & Any cloud option must meet FedRAMP Moderate-equivalent and incident-reporting terms. \\
  Cryptography & FIPS 140-2/140-3 validated modules & FIPS-mode OS and encrypted disks. Confirm that the chosen hypervisor/NAS operating system supports this; an appliance NAS OS may not. \\
  Identity & CAC/PIV or enterprise MFA; role-based access & OIDC/SAML or Kerberos integration for file share, gateway and consoles. \\
  Audit & Who asked what; which documents were retrieved; who accessed which files & Logs forwarded to the SIEM with retention \PH{period}. \\
  Backup and retention & Contract and records rules & Backups stay inside the authorization boundary; encrypted; restore tested. \\
  Derived artifacts & Fine-tuned weights, embeddings and indexes built from CUI & Treated as CUI: stored in the access-controlled registry and covered by backup and sanitization rules. \\
  Supply chain & Model provenance; hardware sourcing & Approved model list; hash-verified offline import; TAA/Section~889-compliant hardware (see Section~\ref{sec:procure}). \\
  Export & ITAR/EAR data may be present & On-premises avoids foreign-hosted data; U.S.-person access control. \\
  Network & Air-gapped or enclave-only & Offline import of models and packages; internal mirrors. \\
  Authority & ATO / System Security Plan & One boundary should cover both computers; plan the authorization effort. \\
  \bottomrule
\end{xltabular}}

\ReportSubsection{Key assumptions (to validate)}
\TodoAnchor{aoa-assump}
\begin{enumerate}
  \item At least 20 named LLM users, with about 12 using the model at the same time
    (both stated by the sponsor). That is about 60\% simultaneous use, far above the
    10--20\% typical of office use, so the design is sized to 12 simultaneous users
    rather than to a ratio. Expected growth in named users: \PH{N}.
  \item Typical prompt 4--16K tokens (with retrieval); long jobs up to 128K.
    Typical response 300--1,500 tokens.
  \item Targets: time to first token under 3~s (short prompt) and under 20~s
    (128K prompt); at least 15--20 tokens/s per user while streaming.
  \item ``Training'' means fine-tuning and adaptation of existing open-weight
    models and small custom models. Pre-training a new LLM is out of scope at
    this budget.
  \item Failover target: \PH{RTO, e.g., manual switchover within 1~hour} for
    LLM serving. The file share and license services have their own recovery
    targets: \PH{RTO/RPO}.
  \item A controlled machine room or rack with 208/240~V circuits is available
    (GPU nodes are loud and hot; an open office is not suitable).
  \item Open-weight models are permitted after security review.
  \item Hardware is ordered from Dell, configured to specification. Prices here are street
    prices with a planning multiplier of 1.3 (range 1.15--1.4) until a quote is received.
  \item Storage need (best estimate, to confirm): the archive is about 20--40~TB now and
    grows 20--25\% per year, so Computer~1 needs about 40--70~TB usable within 3 years,
    including snapshots, vector indexes and model artifacts.
  \item The labor assumptions in Section~\ref{sec:tco} were accepted by the sponsor for now.
\end{enumerate}
